\chapter{Conclusions and recommendations}
\label{chap:conclusion}

This chapter states what the project achieved against each objective, what
the results mean for the design of these details, and what should be done
next.

\section{Conclusions}
\label{sec:conclusions}

The project set out to build and verify a desktop tool that sizes three
marine structural details from a plain description of the job, checks them
against published rules, confirms them with a built-in finite element
analysis and issues a calculation package that a non-programmer can produce
and review. \NEMO{}~2 does this for lifting padeyes, pipe support brackets and
stabilizer fins. The objectives were met as follows.

\begin{enumerate}
    \item \textbf{Rule-based check models.} The padeye is checked to ASME
    BTH-1, the bracket to AISC 360-16 in allowable strength design with a
    check of the pipe span against five machinery excitations, and the fin to
    IACS UR S10 with a Whicker-Fehlner foil load and a Hess-Smith panel method.
    Every check returns its working and names its source, and the 15 padeye,
    14 bracket and 11 fin checks separate code rules from practice. Objective
    1 was met.
    \item \textbf{Stock-aware sizing.} The search enumerates the plate stock,
    shackles, U-bolts and weld sizes, sizes the continuous dimensions with a
    bounded Nelder-Mead search, rounds to buildable sizes and returns designs
    at 100\%, 85\% and 70\% maximum use. Lower-bound pruning skipped 60 of 63
    combinations in the bracket case. Objective 2 was met.
    \item \textbf{Open-source geometry and analysis.} Every design is built as
    an exact solid, meshed three times and solved with quadratic tetrahedra.
    Force balance closed to $10^{-11}$ or better and the assessed peak changed
    by 0.54\% or less between the two finest meshes in all three case studies.
    The manual stress step that blocked the first phase is gone. Objective 3
    was met.
    \item \textbf{Verification.} The engine matches a padeye hand calculation
    within 0.002\%, CAD volumes exactly, beam theory within 2.6\%, a published
    stress concentration factor within 0.9\% and foil theory within 1.7\%. All
    174 automated tests pass. Objective 4 was met for code verification. A
    comparison with a published BTH-1 worked example is still to be done.
    \item \textbf{Interface and outputs.} The tool installs from one file and
    runs in the browser. It shows the engineering steps live, gives the
    working of every check, and writes a PDF package, STEP, STL and DXF files.
    The Fusion connector builds a parametric padeye and re-checks edited
    dimensions. Objective 5 was met in build. In its first Fusion session the
    connector built a padeye with the same volume as \NEMO{}'s, but no agreeing
    stress cross-check has been made yet, and ease of use has not been measured with
    users.
    \item \textbf{Case studies.} One case per part was run, with sweeps of
    load, sling angle, pipe size, support spacing and ship speed, and the
    limits of the results were stated in Section~\ref{sec:limitations}.
    Objective 6 was met.
\end{enumerate}

The case studies also gave findings about the parts themselves.

\begin{itemize}
    \item A 5~t two-point padeye weighs 0.305~kg and is governed by its base
    weld at 98\%. Asking for 15\% spare capacity raises the mass by 42\%,
    because the shackle has to move up a size. Catalogue steps, more than plate
    dimensions, set the mass of small lifting details.
    \item An NPS~4 seawater line bracket uses 25\% of its strength and is
    sized by minimum plate, weld and U-bolt sizes. Its pipe span limits support
    spacing to 2.3~m because of engine firing at 37.5~Hz. Shortening the 3~m
    span to 2.5~m would move it into resonance, so pass or fail does not
    improve steadily as spacing falls.
    \item An oblique sling at 45$^\circ$ to the plate of the 5~t padeye cuts the
    shackle's rating to 75\% and raises the mass by 41\%. For large pipes the
    bracket's own flexibility lowers the span frequency by up to a quarter.
    Close stock bearings can make the fin stock larger than the stress rules
    alone would.
    \item For a 1.4~m by 1.2~m fin at 18~knots, the foil model gives a force
    21\% above the UR S10 rudder formula. The stock governs at 94.5\%, the
    plating sits at rule minimums, and leading-edge cavitation begins between
    18 and 20~knots at the largest angle.
\end{itemize}

The common lesson is that a design method for these details cannot stop at
strength. Minimum sizes, catalogue steps, pipe vibration and foil
hydrodynamics decided the design in every case study.

\section{Recommendations}
\label{sec:recommendations}

For engineers using \NEMO{}~2:

\begin{itemize}
    \item Use the calculation package as a checked starting point, and have a
    competent engineer review it before fabrication. The package is written so
    that the review can be done without the software.
    \item Enter ship-motion accelerations from the classification society
    rules for the vessel, not the defaults, when designing brackets.
    \item Treat the recommended support spacing as the limit for pipe runs
    near machinery, even where the actual spacing passes the separation check.
    \item For fins, use the larger of the two load methods, as the tool does,
    and treat a cavitation warning as a limit on the largest fin angle at that
    speed.
    \item Run the Verification page after installing the tool on a new
    computer, and keep its result with the project file.
\end{itemize}

For the department:

\begin{itemize}
    \item The tool could be used in teaching structural design of ship
    details, because every check shows its working and the effect of an input
    can be seen in a few minutes.
    \item Future projects could add parts to the same engine, such as
    foundations for small machinery or deck fittings, using the existing check
    record, rule sets, sizing method and analysis pipeline.
\end{itemize}

\section{Future work}
\label{sec:future-work}

\begin{enumerate}
    \item Compare the padeye checks with a published ASME BTH-1 worked
    example.
    \item Repeat the Fusion Static Stress cross-check of a padeye with
    connector version 1.1.0, which sets the material in Fusion, and record an
    agreeing result. Extend the connector to brackets and fins.
    \item Carry out a usability study in which marine engineering students who
    have not used the tool complete a design from job description to
    calculation package, and record their time, errors and comments.
    \item Proof load test a padeye designed by the tool, with strain gauges
    near the hole, to validate the finite element model.
    \item Compute bracket accelerations from a classification society method,
    such as the DNV rules for ship motions, from the vessel's main dimensions.
    \item Extend the pipe span model to continuous runs with bends, valves and
    flexible supports.
    \item Add dynamic and fatigue checks for the fin, and a model of the whole
    fin, not only the stock, in the finite element analysis.
    \item Add rule sets from classification societies, such as DNV, Lloyd's
    Register and Bureau Veritas, through the existing rule-set interface.
\end{enumerate}
